
\documentclass{article}

\usepackage{amsmath,amsthm}

\newtheorem{corollary}{Corollary}

\newcommand{\WSPT}{\mathrm{WSPT}}
\newcommand{\OPT}{\mathrm{OPT}}
\newcommand{\Cmax}{C_{\max}}
\newcommand{\Cmin}{C_{\min}}

\begin{document}
\begin{corollary}\label{Corr: Splitting Jobs}
For every instance $I$ of the $P||\sum_j p_j C_j(\alpha)$ problem with $\alpha \in [\frac12,1]$, there is an
instance $I'$ of the $P||\sum_j p_j C_j(\alpha)$ problem such that
$\lambda(I')\geq\lambda(I)$, and every job in $I'$ that starts before
$C_{max}^{min}(S_{\WSPT}(I'))$ has arbitrarily small processing time.
\end{corollary}

\begin{proof}
Assume an arbitrary list schedule $S$ for an instance $I$ of
$P||\sum_j p_j C_j (\alpha)$. We generate an instance $J$ by taking a job $j_0$ from
$I$ and replacing it in schedule $\sigma(I)$ with two jobs $j_1$ and $j_2$ such
that. (kind of inexact to call both the schedule for J and for I S....)
\[
        p_{j_1}+p_{j_2}=p_{j_0},
        \qquad
        C_{j_2}=C_{j_0},
        \qquad
        C_{j_1}=C_{j_0}-p_{j_2}.
\]


This replacement does not alter the completion time of any other job in
$\sigma(I)$. Therefore,
\[
        C_{\max}^i(\sigma(I))=C_{\max}^i(S(J))
\]
for all $1\leq i\leq m$. By  (Queyraynne ref. here)
\begin{align*}
        \sum_{j\in I}p_j C_j(\sigma(I))(\alpha)-\sum_{j\in J}p_j C_j(S(J))(\alpha) 
        &=
        \frac12((2\alpha-1)(p_{j_0}^2-p_{j_1}^2-p_{j_2}^2))\\
        &=
        (2\alpha-1)p_{j_1}p_{j_2}\\
        &\ge 0.
\end{align*}
If $C_{j_1}\leq C_{max}^{min}(S_{\WSPT}(I))$, then the resulting schedule
$S_{\WSPT}(J)$ is again a list schedule. Since every job order is an WSPT order
for an instance of $P||\sum_j p_j C_j(\alpha)$, we may choose an order for $J$ such
that $S(J)=S_{\WSPT}(J)$.

Because $S_{\OPT}(J)$ is also a list schedule, we obtain
\[
        1\leq \lambda(I)
        =
        \frac{C(I)}{C^*(I)}
        \leq
        \frac{C(I)-(2\alpha-1)p_{j_1}p_{j_2}}{C^*(I)-(2\alpha-1)p_{j_1}p_{j_2}}
        \leq
        \frac{C(J)}{C^*(J)}
        =
        \lambda(J).
\]
Repeating this splitting operation yields an instance in which all jobs that
start before $C_{max}^{min}(S_{\WSPT}(I))$ have arbitrarily small processing time.
Since at most $m-1$ jobs can start at $C_{max}^{min}(S_{\WSPT}(I))$, all but at most
$m-1$ jobs have arbitrarily small processing time.
\end{proof}

By the preceding reductions, we from now on only consider instances $I$
with $w_j=p_j$ for all jobs and in which all jobs have arbitrarily small
processing time, except for at most $m-1$ long jobs that start at
$C_{max}^{min}(S_{\WSPT}(I))$. Let $L(I)\subseteq I$ denote the set of these long jobs.

Furthermore, let $K(I)\subseteq L(I)$ be the largest set such that
\[
        p_j\geq \frac1{m-|K(I)|}\sum_{j\in I\setminus K(I)}p_j
\]
holds for all $j\in K(I)$. We obtain $K(I)$ by adding long jobs in decreasing
order of their processing times and testing this property. All jobs in
$I\setminus L(I)$ are called small jobs.

Since every small job has arbitrarily small processing time, we can ignore
the term $(2\alpha-1)\sum_{j\in I\setminus L(I)}p_j^2$ in
(Queryanne ref. here or label of Corr. 2). Thus, for any list schedule $\sigma(I)$,
\begin{equation}
\label{eq:long-job-identity}
        \sum_{j\in I}p_j C_j
        =
        \frac12
        \left(
            (2\alpha-1)\sum_{j\in L(I)}p_j^2
            +
            \sum_{i=1}^m(C_{\max}^i(\sigma(I)))^2
        \right).
\end{equation}
Applying \eqref{eq:long-job-identity} to $S_{\WSPT}(I)$ gives
\begin{equation}
\label{eq:wspt-value}
\begin{aligned}
        C(I)
        &=
        \frac12(2\alpha-1)\sum_{j\in L(I)}p_j^2
        +
        \frac12\sum_{j\in L(I)}
            (p_j+C_{max}^{min}(S_{\WSPT}(I)))^2                    \\
        &\quad
        +
        \frac12\left((m-|L(I)|)(C_{max}^{min}(S_{\WSPT}(I)))^2\right).
\end{aligned}
\end{equation}

Next, we address the construction of an optimal schedule for such instances.
Consider a job order for $I$ in which the long jobs precede all small jobs.
In the corresponding list schedule $S_{\OPT}(I)$, all long jobs start at time
$0$, because $|L(I)|<m$. We have,
\[
        C_{max}^{min}(S_{\OPT}(I))
        \leq
        \frac1{m-|K(I)|}\sum_{j\in I\setminus K(I)}p_j .
\]
If a small job $j_0$ finishes after
\[
        \frac1{m-|K(I)|}\sum_{j\in I\setminus K(I)}p_j ,
\]
we split $j_0$ at time
\[
        \frac1{m-|K(I)|}\sum_{j\in I\setminus K(I)}p_j
\]
as in the proof of \ref{Corr: Splitting Jobs}. In the list for $S_{\OPT}(I)$, the first new job
$j_1$ replaces $j_0$, while the second new job $j_2$ is placed at the end of
the list. After at most $m-|K(I)|-1$ such splitting steps,
\[
        C_{\max}^i(S_{\OPT}(I))
        =
        C_{max}^{min}(S_{\OPT}(I))
        =
        \frac1{m-|K(I)|}\sum_{j\in I\setminus K(I)}p_j
\]
holds for every machine $i$ that executes a small job. Therefore, every
machine $i$ with
\[
        C_{\max}^i(S_{\OPT}(I))>C_{max}^{min}(S_{\OPT}(I))
\]
executes only a single long job. By \eqref{eq:long-job-identity} and the
convexity of the square function, it is impossible to further reduce
\[
        \sum_{i=1}^m(C_{\max}^i(S_{\OPT}(I)))^2
\]
by transferring jobs. Hence, $S_{\OPT}(I)$ is optimal.


\end{document}
